% Part: normal-modal-logic
% Chapter: filtrations
% Section: euclidean-filtrations

\documentclass[../../../include/open-logic-section]{subfiles}

\begin{document}

\olfileid{nml}{fil}{euc}

\olsection{Filtrasi Model Euklidean}

Cara ing \olref[acc]{sec} ora lumaku kanggo model kang Euklidean
utawa serial lan Euklidean. Gatekna model ing pérangan ndhuwur
\olref{fig:ser-eucl}, kang Euklidean lan serial. Upamane
$\Gamma = \{p, \Box p \}$. Nalika nggawe filtrasi liwat $\Gamma$,
mula $[w_1] = [w_3]$ amarga $w_1$ lan $w_3$ iku siji-sijine jagad
kang padha nilaine tumrap~$\Gamma$. Saben filtrasi uga nduweni panah
warisan saka $\mModel{M}$, kaya ing \olref{fig:ser-eucl2}. Model iku
ora Euklidean. Kajaba iku, kita ora bisa nambah panah ing model iku
supaya dadi Euklidean. Kita kudu nambah panah bolak-balik antarane
$[w_2]$ lan $[w_4]$, banjur uga antarane $[w_2]$ lan~$[w_5]$. Nanging $\Box
p$ kudune bener ing~$w_2$, dene $p$ salah ing~$w_5$.

\begin{figure}[htpb]
  \centering
  \begin{tikzpicture}[modal]
    \node[world] (w1)
          [label={left: \mFalse{p}},
            label={below: $\mSat{{}}{\Box p}$}] {$w_1$}; 
    \node[world] (w2)
          [label={right: \mTrue{p}},
            label = {below: $\mSat{{}}{\Box p}$},
            right=of w1] {$w_2$}; 
    \draw[->] (w1) to (w2);
    \draw[reflexive above] (w2) to (w2);
    \node[world] (w3)
          [label={left: \mFalse{p}},
            label={below: $\mSat{{}}{\Box p}$},
            below=of w1] {$w_3$}; 
    \node[world] (w4)
          [label={right:\mTrue{p}},
            label={below: $\mSat/{{}}{\Box p}$},
            right=of w3] {$w_4$};
    \draw[->] (w3) to (w4);
    \draw[reflexive above] (w4) to (w4);     
    \node[world] (w5)
          [label={right: \mFalse{p}},
            label=below:$\mSat/{{}}{\Box p}$,
            right=of w4] {$w_5$}; 
    \draw[->, bend left] (w4) to (w5);
    \draw[->, bend left] (w5) to (w4);
    \draw[reflexive above] (w5) to (w5);     
  \end{tikzpicture}
  \caption{Model serial lan Euklidean.}
  \ollabel{fig:ser-eucl}
\end{figure}

\begin{figure}[ht]
  \centering
  \begin{tikzpicture}[modal,every text node part/.style={align=center}]
    \node[world] (w1)
          [label={left: \mFalse{p}},
            label={right:$[w_1]=[w_3]$},
            label={below: $\mSat{{}}{\Box p}$}] {$[w_1]$}; 
    \node[world] (w2)
         [label={right: \mTrue{p}},
           label = {below: $\mSat{{}}{\Box p}$},
           right=of w1, yshift=1.5cm] {$[w_2]$}; 
    \draw[->] (w1) to (w2);
    \draw[reflexive above] (w2) to (w2);
    \node[world] (w4)
         [label={right:\mTrue{p}},
           label={below: $\mSat/{{}}{\Box p}$},
           right=of w1,yshift=-1.5cm] {$[w_4]$};
    \draw[->] (w1) to (w4);
    \draw[reflexive above] (w4) to (w4);     
    \node[world] (w5)
         [label={right: \mFalse{p}},
             label=below:$\mSat/{{}}{\Box p}$,
             right=of w4] {$[w_5]$}; 
    \draw[->, bend left] (w4) to (w5);
    \draw[->, bend left] (w5) to (w4);
    \draw[reflexive above] (w5) to (w5);     
  \end{tikzpicture}
  \caption{Filtrasi model ing \olref{fig:ser-eucl}.}
  \ollabel{fig:ser-eucl2}
\end{figure}

Mligine, kanggo ngasilake filtrasi Euklidean, ora cukup yen mung
nggatekake filtrasi liwat $\Gamma$ sembarang kang katutup tumrap
sub-!!{formula}s. Kosok baline, kita kudu nggatekake himpunan $\Gamma$
kang \emph{katutup modal} (delengen \olref[pre]{defn:modallyclosed}).
Himpunan ukara mangkono ora winates, mula ora langsung ngasilake
sipat model winates utawa katerputusan kanggo sistem kang cocog.

\begin{thm}\ollabel{thm:modal-closed-filt}
  Upamane $\Gamma$ katutup modal, $\mModel{M}=\tuple{W,R,V}$,
  lan $\mModel{M^*} = \tuple{W^*,R^*,V^*}$ sawijining filtrasi paling
  kasar saka $\mModel{M}$.
  \begin{enumerate}
  \item Yen $\mModel{M}$ simetris, $\mModel{M^*}$ uga simetris.
  \item Yen $\mModel{M}$ transitif, $\mModel{M^*}$ uga transitif.
  \item Yen $\mModel{M}$ Euklidean, $\mModel{M^*}$ uga Euklidean.
  \end{enumerate}
\end{thm}

\begin{proof}
\begin{enumerate}
  \item Latihan. Gunakna kasunyatan yen \Ax{B} lan $\Ax{B_\Diamond}$
    valid ing saben model simetris.
  \item Yen $\mModel{M^*}$ iku filtrasi paling kasar, miturut definisi
    $R^*[u][v]$ kawujud yen lan mung yen $C_1(u,v)$. Kanggo transitivitas,
    upamane $C_1(u,v)$ lan $C_1(v,w)$; kita kudu nuduhake
    $C_1(u,w)$. \iftag{prvBox}{Upamane $\Box!A \in \Gamma$ lan $\mSat{M}{\Box !A}[u]$; mula
      $\mSat{M}{\Box\Box!A}[u]$ amarga \Ax{4} valid ing saben
      model transitif; amarga $\Box\Box!A \in \Gamma$ miturut katutupan,
      $C_1(u,v)$ uga njalari $\mSat{M}{\Box!A}[v]$, lan $C_1(v,w)$
      uga njalari $\mSat{M}{!A}[w]$. }{}\iftag{prvDiamond}{Upamane
      $\Diamond!A \in \Gamma$ lan $\mSat{M}{!A}[w]$; mula $\mSat{M}{\Diamond !A}[v]$ miturut
      $C_1(v,w)$, amarga $\Diamond !A \in \Gamma$ miturut katutupan modal. Miturut
      $C_1(u,v)$, kita entuk $\mSat{M}{\Diamond\Diamond !A}[u]$ amarga
      $\Diamond\Diamond!A \in \Gamma$ miturut katutupan modal. Amarga
      $\Ax{4}_\Diamond$ valid ing saben model transitif,
      $\mSat{M}{\Diamond!A}[u]$.}{}
  \item Latihan. Gunakna kasunyatan yen \Ax{5} lan $\Ax{5_\Diamond}$
    loro-lorone valid ing saben model Euklidean.
\end{enumerate}
\end{proof}

\begin{prob}
  Jangkepna bukti \olref[nml][fil][euc]{thm:modal-closed-filt}.
\end{prob}

\end{document}
